\begin{center}
    {\LARGE \textbf{Parte IV}}\\[6pt]
\end{center}
\vspace{4pt}\hrule\vspace{10pt}

{\large Análisis del Tipo de Interés y la Cuenta Corriente \par}

\vspace{0.1cm}
\noindent\rule{\textwidth}{0.5pt}
\vspace{0.1cm}

\section*{Planteamiento}
Usaremos el resultado de la Parte III para la tasa de interés mundial de equilibrio:
\begin{equation}
1 + \rs = \frac{\sum_{i=1}^{N} \dfrac{y_{i,2}}{1+\beta_i}}{\sum_{i=1}^{N} \dfrac{\beta_i y_{i,1}}{1+\beta_i}}.
\label{eq:tasa_general}
\end{equation}

Ahora dividimos los $N$ países en dos grupos:
\begin{itemize}[leftmargin=1.2cm]
    \item ahorradores $(S)$, con factor de descuento común $\bs$,
    \item endeudados $(B)$, con factor de descuento común $\bb$,
\end{itemize}
con $\bs > \bb$.

Además, denotamos:
\[
\ns = \operatorname{card}(S), \qquad \nb = \operatorname{card}(B), \qquad \ns + \nb = N,
\]
y suponemos dotación simétrica dentro de cada grupo:
\[
(y_{S,1},y_{S,2}) \quad \text{para cada país de } S,
\qquad
(y_{B,1},y_{B,2}) \quad \text{para cada país de } B.
\]

\section*{Ítem 20}

Partiendo de \eqref{eq:tasa_general}, y usando que todos los países dentro de cada grupo son idénticos, el numerador se puede escribir como
\[
\sum_{i=1}^{N} \frac{y_{i,2}}{1+\beta_i}
=
\sum_{i \in S} \frac{y_{S,2}}{1+\bs}
+
\sum_{i \in B} \frac{y_{B,2}}{1+\bb}
=
\ns \frac{y_{S,2}}{1+\bs} + \nb \frac{y_{B,2}}{1+\bb}.
\]

Análogamente, el denominador se puede escribir como
\[
\sum_{i=1}^{N} \frac{\beta_i y_{i,1}}{1+\beta_i}
=
\sum_{i \in S} \frac{\bs y_{S,1}}{1+\bs}
+
\sum_{i \in B} \frac{\bb y_{B,1}}{1+\bb}
=
\ns \frac{\bs y_{S,1}}{1+\bs} + \nb \frac{\bb y_{B,1}}{1+\bb}.
\]

Por lo tanto,
\begin{equation}
1+\rs
=
\frac{\ns \dfrac{y_{S,2}}{1+\bs} + \nb \dfrac{y_{B,2}}{1+\bb}}
{\ns \dfrac{\bs y_{S,1}}{1+\bs} + \nb \dfrac{\bb y_{B,1}}{1+\bb}}.
\label{eq:rgrupos}
\end{equation}

y, en consecuencia,
\begin{equation}
\rs
=
\frac{\ns \dfrac{y_{S,2}}{1+\bs} + \nb \dfrac{y_{B,2}}{1+\bb}}
{\ns \dfrac{\bs y_{S,1}}{1+\bs} + \nb \dfrac{\bb y_{B,1}}{1+\bb}} - 1.
\label{eq:rfinal}
\end{equation}

Para simplificar la notación y simplificar las posteriores derivaciones, definimos
\[
A(\bs,\bb)
=
\ns \frac{y_{S,2}}{1+\bs} + \nb \frac{y_{B,2}}{1+\bb},
\qquad
D(\bs,\bb)
=
\ns \frac{\bs y_{S,1}}{1+\bs} + \nb \frac{\bb y_{B,1}}{1+\bb}.
\]
Entonces
\[
1+\rs = \frac{A}{D},
\qquad
\rs = \frac{A}{D} - 1.
\]

\subsection*{Interpretación en términos de fondos prestables}
La expresión \eqref{eq:rgrupos} puede leerse como un equilibrio entre oferta y demanda de fondos prestables:
\begin{itemize}[leftmargin=1.2cm]
    \item el numerador $A$ recoge la \textbf{oferta agregada de recursos futuros} ponderada por la impaciencia/paciencia de cada grupo;
    \item el denominador $D$ recoge la \textbf{demanda agregada de ahorro presente} ponderada por el factor de descuento de cada grupo.
\end{itemize}

Como el grupo $S$ es más paciente ($\bs > \bb$), ese grupo tiende a ofrecer relativamente más recursos presentes al mercado mundial, mientras que el grupo $B$, al ser menos paciente, tiende a demandarlos.

\subsection*{Derivada de $\rs$ respecto a $\bs$}
Primero derivamos $A$ y $D$ respecto a $\bs$:
\[
\frac{\partial A}{\partial \bs}
=
\frac{\partial}{\partial \bs}
\left(
\ns \frac{y_{S,2}}{1+\bs} + \nb \frac{y_{B,2}}{1+\bb}
\right)
=
-\ns \frac{y_{S,2}}{(1+\bs)^2},
\]

Ahora,
\[
\frac{\partial D}{\partial \bs}
=
\frac{\partial}{\partial \bs}
\left(
\ns \frac{\bs y_{S,1}}{1+\bs} + \nb \frac{\bb y_{B,1}}{1+\bb}
\right)
=
\ns y_{S,1}
\frac{\partial}{\partial \bs}
\left(
\frac{\bs}{1+\bs}
\right).
\]
Usando regla del cociente,
\[
\frac{\partial}{\partial \bs}
\left(
\frac{\bs}{1+\bs}
\right)
=
\frac{(1+\bs)(1) - \bs(1)}{(1+\bs)^2}
=
\frac{1}{(1+\bs)^2}.
\]
Luego,
\[
\frac{\partial D}{\partial \bs}
=
\ns \frac{y_{S,1}}{(1+\bs)^2}.
\]

Como
\[
\rs = \frac{A}{D} - 1,
\]
se tiene
\[
\frac{\partial \rs}{\partial \bs}
=
\frac{\partial}{\partial \bs}\left(\frac{A}{D}\right)
=
\frac{\left(\dfrac{\partial A}{\partial \bs}\right)D - A\left(\dfrac{\partial D}{\partial \bs}\right)}{D^2}.
\]
Sustituyendo,
\[
\frac{\partial \rs}{\partial \bs}
=
\frac{\left(-\ns \dfrac{y_{S,2}}{(1+\bs)^2}\right)D - A\left(\ns \dfrac{y_{S,1}}{(1+\bs)^2}\right)}{D^2}.
\]

El signo es negativo porque $\ns>0$, $(1+\bs)^2>0$, $D^2>0$, $y_{S,2}>0$, $y_{S,1}>0$ y $A,D>0$.

\subsection*{Derivada de $\rs$ respecto a $\bb$}
De forma análoga,
\[
\frac{\partial A}{\partial \bb}
=
-\nb \frac{y_{B,2}}{(1+\bb)^2},
\]
porque solo cambia el término del grupo $B$ en el numerador.

Asimismo,
\[
\frac{\partial D}{\partial \bb}
=
\nb y_{B,1}
\frac{\partial}{\partial \bb}
\left(
\frac{\bb}{1+\bb}
\right)
=
\nb \frac{y_{B,1}}{(1+\bb)^2}.
\]

Entonces,
\[
\frac{\partial \rs}{\partial \bb}
=
\frac{\left(\dfrac{\partial A}{\partial \bb}\right)D - A\left(\dfrac{\partial D}{\partial \bb}\right)}{D^2}
=
\frac{\left(-\nb \dfrac{y_{B,2}}{(1+\bb)^2}\right)D - A\left(\nb \dfrac{y_{B,1}}{(1+\bb)^2}\right)}{D^2}.
\]
Por lo tanto, similar al caso anterior, el impacto es negativo 

\subsection*{¿Por qué el efecto de $\bs$ sobre $\rs$ no es simétrico al efecto de $\bb$?}
Los dos efectos no son simétricos porque las expresiones dependen de parámetros distintos:
\[
\ns,\ y_{S,1},\ y_{S,2},\ \bs
\qquad \text{vs} \qquad
\nb,\ y_{B,1},\ y_{B,2},\ \bb.
\]

Esto se evidencia en las ecuaciones presentadas, en las que la magnitud del efecto depende también del tamaño de cada grupo (ahorradores y endeudados) y de sus dotaciones respectivas .

Para que ambos efectos tengan igual magnitud debería cumplirse:
\begin{equation}
\frac{\ns}{(1+\bs)^2}\left(y_{S,2}D + y_{S,1}A\right)
=
\frac{\nb}{(1+\bb)^2}\left(y_{B,2}D + y_{B,1}A\right).
\label{eq:condicion_general}
\end{equation}

Si las dotaciones por país fueran iguales entre grupos y la única fuente de asimetría es el número de países, entonces la condición a cumplirse sería:
\[
\operatorname{card}(S)=\operatorname{card}(B).
\]

\section*{ítem 21}
En la Parte III obtuvimos, para cualquier país $i$,
\[
CA_{i,1}
=
\frac{\beta_i y_{i,1} - \dfrac{y_{i,2}}{1+\rs}}{1+\beta_i}.
\]

Por lo tanto, para un país representativo del grupo $S$,
\begin{equation}
\cas
=
\frac{\bs y_{S,1} - \dfrac{y_{S,2}}{1+\rs}}{1+\bs},
\label{eq:cas}
\end{equation}
y para un país representativo del grupo $B$,
\begin{equation}
\cab
=
\frac{\bb y_{B,1} - \dfrac{y_{B,2}}{1+\rs}}{1+\bb}.
\label{eq:cab}
\end{equation}

\subsection*{Derivada de $\cas$ respecto a $\bs$}
Tomemos
\[
\cas = \frac{N_S}{1+\bs},
\qquad
N_S = \bs y_{S,1} - \frac{y_{S,2}}{1+\rs}.
\]
Entonces,
\[
\frac{\partial \cas}{\partial \bs}
=
\frac{(1+\bs)\dfrac{\partial N_S}{\partial \bs} - N_S}{(1+\bs)^2}.
\]
Ahora,
\[
\frac{\partial N_S}{\partial \bs}
=
\frac{\partial}{\partial \bs}
\left(
\bs y_{S,1} - \frac{y_{S,2}}{1+\rs}
\right)
=
 y_{S,1} + \frac{y_{S,2}}{(1+\rs)^2}\frac{\partial \rs}{\partial \bs},
\]
porque
\[
\frac{\partial}{\partial \bs}\left(-\frac{y_{S,2}}{1+\rs}\right)
=
-y_{S,2}\left(-\frac{1}{(1+\rs)^2}\right)\frac{\partial \rs}{\partial \bs}
=
\frac{y_{S,2}}{(1+\rs)^2}\frac{\partial \rs}{\partial \bs}.
\]

Sustituyendo en la regla del cociente y simplificando,
\begin{equation}
\frac{\partial \cas}{\partial \bs}
=
\frac{y_{S,1}}{(1+\bs)^2}
-
\frac{y_{S,2}((1+\rs)+(1+\bs))}{(1+\bs)^2(1+\rs)^2}\frac{\partial \rs}{\partial \bs}.
\label{eq:dcasdbetas}
\end{equation}

Como $\dfrac{\partial \rs}{\partial \bs}<0$, el segundo término es positivo debido al signo negativo por delante. De este modo, junto al primer término, obtenemos un impacto general positivo:
\[
\frac{\partial \cas}{\partial \bs} > 0.
\]

\subsection*{Derivada de $\cas$ respecto a $\bb$}
Aquí $\bs$ no cambia directamente, así que solo opera el efecto indirecto vía $\rs$:
\begin{equation}
\frac{\partial \cas}{\partial \bb}
=
\frac{y_{S,2}}{(1+\bs)(1+\rs)^2}\frac{\partial \rs}{\partial \bb}.
\label{eq:dcasdbetab}
\end{equation}
Como $\dfrac{\partial \rs}{\partial \bb}<0$, se concluye que
\[
\frac{\partial \cas}{\partial \bb}<0.
\]

\subsection*{Derivada de $\cab$ respecto a $\bs$}
De forma análoga,
\begin{equation}
\frac{\partial \cab}{\partial \bs}
=
\frac{y_{B,2}}{(1+\bb)(1+\rs)^2}\frac{\partial \rs}{\partial \bs}.
\label{eq:dcabdbetas}
\end{equation}
Como $\dfrac{\partial \rs}{\partial \bs}<0$, entonces
\[
\frac{\partial \cab}{\partial \bs}<0.
\]
Esto significa que un aumento de $\bs$ vuelve más deficitario al grupo $B$.

\subsection*{Derivada de $\cab$ respecto a $\bb$}
Ahora sí hay un efecto directo positivo sobre la cuenta corriente del propio grupo, además del efecto indirecto vía $\rs$:
\begin{equation}
\frac{\partial \cab}{\partial \bb}
=
\frac{y_{B,1}}{(1+\bb)^2}
-
\frac{y_{B,2}((1+\rs)+(1+\bb))}{(1+\bb)^2(1+\rs)^2}\frac{\partial \rs}{\partial \bb}.
\label{eq:dcabdbetab}
\end{equation}
Como en el caso anterior, finalmente obtenemos un impacto positivo:
\[
\frac{\partial \cab}{\partial \bb}>0.
\]

\subsection*{Conclusión sobre los signos}
En resumen,
\[
\frac{\partial \cas}{\partial \bs} > 0,
\qquad
\frac{\partial \cab}{\partial \bs} < 0,
\qquad
\frac{\partial \cas}{\partial \bb} < 0,
\qquad
\frac{\partial \cab}{\partial \bb} > 0.
\]

Por tanto:
\begin{itemize}[leftmargin=1.2cm]
    \item un aumento en $\bs$ genera un \textbf{mayor superávit} en el grupo $S$ y un \textbf{mayor déficit} en el grupo $B$;
    \item un aumento en $\bb$ produce el efecto contrario: reduce el superávit de $S$ y mejora la cuenta corriente de $B$.
\end{itemize}

\subsection*{Interpretación de la matriz}
La matriz
\[
\begin{pmatrix}
\dfrac{\partial \cas}{\partial \bs} & \dfrac{\partial \cas}{\partial \bb} \\
\dfrac{\partial \cab}{\partial \bs} & \dfrac{\partial \cab}{\partial \bb}
\end{pmatrix}
\]
puede interpretarse como una \textbf{matriz de sensibilidad de las cuentas corrientes} con respecto a los factores de descuento de ambos grupos.

Su interpretación económica es la siguiente:
\begin{itemize}[leftmargin=1.2cm]
    \item los elementos de la diagonal principal son positivos: un aumento en la paciencia de un grupo mejora su propia cuenta corriente;
    \item los elementos fuera de la diagonal son negativos: un aumento en la paciencia de un grupo empeora la cuenta corriente del otro grupo;
    \item la matriz resume cómo cambios en preferencias intertemporales de cada tipo de país afectan el ahorro y el endeudamiento dentro de la economía mundial.
\end{itemize}




